\documentclass[a4paper]{tufte-book}

\title{The Parachute Book, Vol I: Bare Metal}
\author{ Matt Gumbley, DevZendo.org }

%\date % without \date command, current date is supplied

%\geometry{showframe} % display margins for debugging page layout


% The tufte-latex package is great for pdflatex, but TeXPad uses xelatex, which fouls up the header, and doesn't build unless the following section is added. The header doesn't use a smallcaps font, and I had to change the main text and math font to make it look a bit better.

% Thanks to https://tex.stackexchange.com/questions/200722/xetex-seems-to-break-headers-in-tufte-handout
\usepackage{fontspec}
\usepackage{ifxetex}


\ifxetex
  \newcommand{\textls}[2][5]{%
    \begingroup\addfontfeatures{LetterSpace=#1}#2\endgroup
  }
  \renewcommand{\allcapsspacing}[1]{\textls[15]{#1}}
  \renewcommand{\smallcapsspacing}[1]{\textls[10]{#1}}
  \renewcommand{\allcaps}[1]{\textls[15]{\MakeTextUppercase{#1}}}
  \renewcommand{\smallcaps}[1]{\smallcapsspacing{\scshape\MakeTextLowercase{#1}}}
  \renewcommand{\textsc}[1]{\smallcapsspacing{\textsmallcaps{#1}}}
\fi

\setmainfont[
  Renderer=Basic,
  Numbers=OldStyle,
  Scale=1.0,
  Extension=.otf,
  UprightFont=*-regular,
  ItalicFont=*-italic,
  BoldFont=*-bold,
  BoldItalicFont=*-bolditalic]{texgyrepagella} %% The Palatino from the TeX Gyre Project

% Get the small caps right. Too much fiddling going on when using xelatex and tufte-latex.
% The title and author aren't as nice as with pdflatex, but at least it builds.
% Path = /usr/share/texmf/fonts/opentype/public/tex-gyre/ ,
\setmainfont{texgyrepagella-regular.otf}[
  BoldFont       = texgyrepagella-bold.otf,
  ItalicFont     = texgyrepagella-italic.otf,
  BoldItalicFont = texgyrepagella-bolditalic.otf,
  SmallCapsFeatures={Letters=SmallCaps},
]

% \setsansfont[Renderer=Basic, Scale=0.90]{Helvetica}
% Set up the spacing using fontspec features
\renewcommand\allcapsspacing[1]{{\addfontfeature{LetterSpace=15}#1}}
\renewcommand\smallcapsspacing[1]{{\addfontfeature{LetterSpace=10}#1}}


\usepackage{graphicx} % allow embedded images
\setkeys{Gin}{width=\linewidth,totalheight=\textheight,keepaspectratio}
\graphicspath{{graphics/}} % set of paths to search for images
\usepackage{amsmath}  % extended mathematics
\usepackage{booktabs} % book-quality tables
\usepackage{units}    % non-stacked fractions and better unit spacing
\usepackage{multicol} % multiple column layout facilities
\usepackage{lipsum}   % filler text
\usepackage{fancyvrb} % extended verbatim environments
\usepackage{hyperref} % URLs
\usepackage{xcolor}   % For coloured text
\usepackage{soul}     % For highlighted text
\fvset{fontsize=\normalsize}% default font size for fancy-verbatim environments

% Standardize command font styles and environments
\newcommand{\doccmd}[1]{\texttt{\textbackslash#1}}% command name -- adds backslash automatically
\newcommand{\docopt}[1]{\ensuremath{\langle}\textrm{\textit{#1}}\ensuremath{\rangle}}% optional command argument
\newcommand{\docarg}[1]{\textrm{\textit{#1}}}% (required) command argument
\newcommand{\docenv}[1]{\textsf{#1}}% environment name
\newcommand{\docpkg}[1]{\texttt{#1}}% package name
\newcommand{\doccls}[1]{\texttt{#1}}% document class name
\newcommand{\docclsopt}[1]{\texttt{#1}}% document class option name
\newenvironment{docspec}{\begin{quote}\noindent}{\end{quote}}% command specification environment


\newcommand{\floor}[1]{\left\lfloor #1 \right\rfloor}
\newcommand{\ceil}[1]{\left\lceil #1 \right\rceil}

% Generates the index
\usepackage{makeidx}
\makeindex

%%%%%%%%%%%%%%%%%%%%%%%%%%%%%%%%%%%%%%%%%%%%%%%%%%%%%%%%%%%%%%%%%%%%%%%%%%%%%%%%%%%%%%%%%

\begin{document}

\frontmatter

    \maketitle % this prints the handout title, author, and date

%%%%%%%%%%%%%%%%%%%%%%%%%%%%%%%%%%%%%%%%%%%%%%%%%%%%%%%%%%%%%%%%%%%%%%%%%%%%%%%%%%%%%%%%%



\begin{fullwidth}
~\vfill
\thispagestyle{empty}
\setlength{\parindent}{0pt}
\setlength{\parskip}{\baselineskip}
Copyright \copyright\ \the\year\ \thanklessauthor

% about DevZendo.org ?


\par\smallcaps{https://devzendo.github.io/parachute}

\par Licensed under the Apache License, Version 2.0 (the ``License''); you may not
use this file except in compliance with the License. You may obtain a copy
of the License at \url{http://www.apache.org/licenses/LICENSE-2.0}. Unless
required by applicable law or agreed to in writing, software distributed
under the License is distributed on an \smallcaps{``AS IS'' BASIS, WITHOUT
WARRANTIES OR CONDITIONS OF ANY KIND}, either express or implied. See the
License for the specific language governing permissions and limitations
under the License.\index{license}

\par\textit{Parachute 0.0.2, June 2026}
\end{fullwidth}


%%%%%%%%%%%%%%%%%%%%%%%%%%%%%%%%%%%%%%%%%%%%%%%%%%%%%%%%%%%%%%%%%%%%%%%%%%%%%%%%%%%%%%%%%

    \pagebreak
    \tableofcontents

%%%%%%%%%%%%%%%%%%%%%%%%%%%%%%%%%%%%%%%%%%%%%%%%%%%%%%%%%%%%%%%%%%%%%%%%%%%%%%%%%%%%%%%%%
\pagebreak
~\vfill
\begin{doublespace}
\noindent\fontsize{18}{22}\selectfont\itshape
\nohyphenation

\noindent
In memoriam\\
\vfill
Sir Charles Antony Richard Hoare,\\
FRS, FREng.\\
\vfill
11 January 1934 - 5 March 2026


\end{doublespace}
\vfill
\vfill
%%%%%%%%%%%%%%%%%%%%%%%%%%%%%%%%%%%%%%%%%%%%%%%%%%%%%%%%%%%%%%%%%%%%%%%%%%%%%%%%%%%%%%%%%


% each chapter opens with a brief paragraph stating its scope and purpose, and closes with a summary and forward pointer.

    
\pagebreak
\part{Preface}
\section{About this book}
    \newthought{Parachute is...} an open, simple and useful personal computing stack — built around the Inmos Transputer (in emulation).
        
	Parachute is a long-term project to build a small systems architecture
        from the ground up: emulator, toolchain, OS kernel, and applications — designed
        to run on cheap, widely available hardware, and small enough that a single
        person can understand the whole system.
        
	This book is a tutorial introduction and complete reference to part of the Parachute system: its emulator, hardware, languages, and running small parallel applications on it.

% who this is for, what you need, conventions, how chapters relate, how to get the source code
% Who is reading a document?
% What do they want to accomplish?
% Why do they want to accomplish it?


\subsection{Subsection One}

    Some things...

\subsection{Subsection Two}

    Some more things....
    
\subsection{Typographic conventions}

%%%%%%%%%%%%%%%%%%%%%%%%%%%%%%%%%%%%%%%%%%%%%%%%%%%%%%%%%%%%%%%%%%%%%%%%%%%%%%%%%%%%%%%%%



\pagebreak
\part{Foundations}

% vision/motivation, what a Transputer is, the architecture, the ecosystem (IServer etc.)


Parachute is guided by the ideas of people who believed software could be lean, comprehensible, and purposeful:
\begin{itemize}
	\item The Inmos Transputer, occam, and C.A.R. Hoare's CSP — the technical foundation (with other languages)
	\item Niklaus Wirth's Project Oberon and his essay "A Plea for Lean Software" — a philosophy of restraint
	\item The Amiga OS — elegance and capability in a small footprint
	\item Psion PDAs and the SIBO/EPOC operating systems — productive, focused, astonishingly efficient
	\item The Palm Pilot — proof that constraints produce clarity
\end{itemize}

Think: the early Mac, the Amiga, early Windows — but parallel, distributed, and built on a CPU that was genuinely ahead of its time.
	
And remember the PDA? Pocket-sized. Days of battery life on AAs. Distraction-free and productive. The smartphone was not its natural successor. Touch screens are necessary, but not sufficient for productive work - proper keyboards are essential. 



\section{Project Website}



%%%%%%%%%%%%%%%%%%%%%%%%%%%%%%%%%%%%%%%%%%%%%%%%%%%%%%%%%%%%%%%%%%%%%%%%%%%%%%%%%%%%%%%%%

\pagebreak
\part{Getting Started}

% In the tutorial part, resist the temptation to front-load explanation. The reader should do something that works in the first chapter before you explain why it works. Explanation belongs in Part I or in clearly marked 'under the hood' sections, not interspersed in tutorial steps. This is where many self-authored technical books drift.

\section{Project Website}

The website for the parachute software is located at
\url{https://devzendo.github.io/parachute}. On the website, you'll find
links to our \smallcaps{github} repository, mailing lists, bug tracker, and documentation. 

Releases of parachute may be found in the \smallcaps{releases} area of the \smallcaps{github} repository: \url{https://github.com/devzendo/parachute/releases} .



\section{Downloading}




\section{Verifying downloads}




\section{Installing}
\subsection{Installing on Windows}
\subsection{Installing on macOS}
\subsection{Installing on Linux}
\subsection{Installing on Raspberry Pi Pico}





\pagebreak
\section{Tutorial introduction}
\subsection{Hello World, prebuilt}
\subsection{Assembling Hello World}
\subsection{Building Small-C Hello World}
\subsection{eForth}
\subsection{The Mandlebrot Set}



%%%%%%%%%%%%%%%%%%%%%%%%%%%%%%%%%%%%%%%%%%%%%%%%%%%%%%%%%%%%%%%%%%%%%%%%%%%%%%%%%%%%%%%%%


\pagebreak
\part{Programming}

% tutorials in escalating complexity, one concept per chapter, each self-contained but building on prior ones; multiple languages means likely a chapter per language introducing it before the complexity builds up

%%%%%%%%%%%%%%%%%%%%%%%%%%%%%%%%%%%%%%%%%%%%%%%%%%%%%%%%%%%%%%%%%%%%%%%%%%%%%%%%%%%%%%%%%



\pagebreak
\part{Reference}
% one chapter per tool/language, dense, structured consistently (syntax, semantics, options, error conditions)

\section{Software}
\subsection{Emulator}
\subsection{IServer}
\subsection{eForth}

\section{Languages}
\subsection{Macro Assembler}
\subsection{Small-C Compiler}
\subsection{Small-C Assembler}
\subsection{eForth}



%%%%%%%%%%%%%%%%%%%%%%%%%%%%%%%%%%%%%%%%%%%%%%%%%%%%%%%%%%%%%%%%%%%%%%%%%%%%%%%%%%%%%%%%%



\pagebreak
\appendix
\section*{Appendices}
\addcontentsline{toc}{section}{Appendices}
\renewcommand{\thesubsection}{\Alph{subsection}}

% quick reference tables, platform-specific notes, bibliography


\subsection{Emulator Compatibility}

\subsection{IServer Compatibility}



%%%%%%%%%%%%%%%%%%%%%%%%%%%%%%%%%%%%%%%%%%%%%%%%%%%%%%%%%%%%%%%%%%%%%%%%%%%%%%%%%%%%%%%%%

% The back matter contains appendices, bibliographies, indices, glossaries, etc.

\backmatter


\pagebreak
\bibliography{parachute-book}
\bibliographystyle{plainnat}

\printindex

\end{document}
